\section{}

\paragraph{Overview:}

This week was all about taking my initial learnings and ensuring that my theoretical knowledge would be applicable to a functional artifact. I aimed to have a functioning node mapping system built (the actual note content systems would be built in the coming weeks)

\paragraph{Achievements, Thoughts and Ideas:}
- I decided to take my technical research thus far and produce more demos to create a prototype for the automated and flexible mind map technology I wanted to build. 
- I started the week by cleaning up my prior notes and diary entries to ensure any future progress stays entirely relevant to my project targets and that my time is spent most efficiently
- I explored some preexisting software solutions for notetaking to understand the core features I would need to learn how to build
- I researched two core technologies I could use in my project: Flask to host my website so that users can access it and Cystoscope for a framework to create my graph with
- I then moved on to the creation of 6 prototypes, building new features and improving on my previous ideas 

I watched a video on student self discipline \cite{mikedee13WaysBe2021} and made these notes around motivation and student disciplined:
\begin{enumerate}
	\item Using the 5 second window to move from idea to action. Maximising efficacy and minimising time spent off task
	\item Start small with change, it is much easier to implement many small changes than one large shift
	\item Power of three rule: making a list of three tasks every day that you want to achieve. This idea really spoke to me and I decided to do this each day with regard to my EPQ, inside each diary entry I have now put a power of three rule for each of the days I work on /research for my EPQ. 
	\item Set boundaries regarding social media use to reduce short term gratification in favour of long term success. I already set myself time slots of no distractions (do not disturb on devices and phone left out of study room) but it was nice to learn more.
	\item Visualise future achievement, when feeling vulnerable to procrastination imagine the future success you will see from the productive time spent now
	\item Schedule your down time around your study sessions, and enjoy this time without the distraction of work
	\item "your start does not determine now you finish“ you can always fix your past mistakes with work now
	\item Organize your mind, use note taking methods to help reduce the millions of thoughts that bounce around your mind. Your brain is a powerful yet limited resource, learning to manage that resource is the key to success 
	\item Eliminate "but" from your vocabulary, find solutions not problems when facing tasks, these solutions often become the most powerful tool we have 
\end{enumerate}

\paragraph{What went well?}

After lots of error fixing and checking the prototypes finally function as id want them to leading to much more exciting feature additions in the future
\paragraph{What did i learn?}
Whilst producing the demos I filled in many knowledge gaps that hadn't been covered by my prior research. Until this point I had confirmed the solutions existed but only now have made advancements into what they look like and how I can build them 

\paragraph{What was challenging?}
Throughout the iterations I completed this week I found many limitations and flaws in prior code that had to be rebuilt continually. Although this was frustrating it has taught me a valuable lesson that I should think ahead to further future expansion of features before implementing an idea in a certain way

\paragraph{Plans for next week:}
With a large amount of my technical feasibility study now complete I want to dive back into the learning methods that underpin my whole EP and learn more about existing research that has been completed

